\chapter{การกระตุ้นเมล็ดพันธุ์และการเจริญเติบโตของพืชด้วยพลาสมา}
\label{chap:ch04}

% =========================================================================
% แผนบริหารการสอนประจำบทที่ 4 (Modular Lesson Plan Boxes - 2-Page Standard)
% =========================================================================

\vspace*{-6pt}
\section*{แผนบริหารการสอนประจำบทที่ 4}
\addcontentsline{toc}{section}{แผนบริหารการสอนประจำบทที่ 4}

% --- Page 1 of Lesson Plan ---
\begin{planbox}{หัวข้อเนื้อหาประจำบท}
\begin{enumerate}
    \item ปัญหาการเสื่อมสภาพของเมล็ดพันธุ์และข้อจำกัดของการเตรียมความพร้อมเมล็ดแบบดั้งเดิม
    \item สัณฐานวิทยาของผิวเมล็ดและกลไกการปรับเปลี่ยนสภาพชอบน้ำด้วยพลาสมา
    \item ทฤษฎีฟิสิกส์ของการเปียก: แบบจำลองมุมสัมผัสของยัง (Young), เวนเซล (Wenzel) และแคสซี-แบกซ์เตอร์ (Cassie-Baxter)
    \item การกัดกร่อนระดับนาโนและการแทรกซึมของน้ำและออกซิเจนผ่านชั้นคิวติเคิล
    \item สรีรวิทยาการงอก: การกระตุ้นเอนไซม์แอลฟา-อะไมเลสและการควบคุมสมดุลฮอร์โมนพืช (GA/ABA)
    \item ดัชนีความแข็งแรงของเมล็ดพันธุ์และการเจริญเติบโตในระยะกล้า
    \item สถาปัตยกรรม Edge AI สำหรับการฟีโนไทป์เมล็ดพันธุ์ปริมาณงานสูงด้วย Multi-Task CNN
\end{enumerate}
\end{planbox}

\begin{planbox}{วัตถุประสงค์เชิงพฤติกรรม}
\begin{enumerate}
    \item อธิบายกลไกทางฟิสิกส์และเคมีในการปรับปรุงผิวเปลือกเมล็ดด้วยพลาสมาความดันบรรยากาศได้
    \item คำนวณพลังงานอิสระที่ผิวและการเปลี่ยนแปลงมุมสัมผัสตามแบบจำลองของยังและเวนเซลได้อย่างถูกต้อง
    \item อธิบายบทบาทของอนุมูลอิสระ RONS ต่อการกระตุ้นการทำงานของเอนไซม์ย่อยแป้งและการสลายฮอร์โมนยับยั้งการงอกได้
    \item คำนวณดัชนีความแข็งแรงของเมล็ดพันธุ์ (Seed Vigor Index) และร้อยละความงอกตามเกณฑ์มาตรฐานสากล ISTA ได้
    \item วิเคราะห์การทำงานของโครงข่ายประสาทเทียมตรวจวัดอัตราการแทงรากแรกเกิดของเมล็ดพันธุ์ได้อย่างเป็นระบบ
\end{enumerate}
\end{planbox}

\newpage

% --- Page 2 of Lesson Plan ---
\begin{planbox}{วิธีสอนและกิจกรรมการเรียนการสอน}
\begin{enumerate}
    \item บรรยายเชิงทฤษฎีเกี่ยวกับฟิสิกส์ของผิวสัมผัสและจลนพลศาสตร์การดูดซับน้ำของเมล็ดพันธุ์
    \item สาธิตการทดลองปรับสภาพผิวเมล็ดพันธุ์ข้าวด้วยเครื่องกำเนิดพลาสมา DBD แบบฟลูอิดไดซ์เบด (Fluidized Bed)
    \item ปฏิบัติการวัดมุมสัมผัสของหยดน้ำ (Water Contact Angle) บนผิวเมล็ดก่อนและหลังการทรีตเมนต์ด้วยโกนิโอมิเตอร์
    \item ปฏิบัติการตรวจวิเคราะห์กิจกรรมของเอนไซม์แอลฟา-อะไมเลสด้วยวิธีสเปกโทรโฟโตเมทรีในห้องปฏิบัติการ
    \item สาธิตการใช้งานระบบ Edge AI ประมวลผลภาพเมล็ดงอกแบบเรียลไทม์บนบอร์ดประมวลผลสมองกลฝังตัว
\end{enumerate}
\end{planbox}

\begin{planbox}{สื่อการเรียนการสอน}
\begin{enumerate}
    \item แหล่งกำเนิดพลาสมา DBD ชนิดฟลูอิดไดซ์เบดสำหรับการทรีตเมล็ดพันธุ์ปริมาณมาก
    \item เครื่องวัดมุมสัมผัสของหยดของเหลวความละเอียดสูง (Contact Angle Goniometer)
    \item ตัวอย่างเมล็ดพันธุ์ข้าวขาวดอกมะลิ 105 (KDML 105) และเมล็ดพันธุ์ข้าวโพดเลี้ยงสัตว์
    \item กล้องสเตอริโอไมโครสโคปและเครื่องสเปกโทรโฟโตมิเตอร์ UV-Vis
    \item บอร์ดประมวลผล Edge AI พร้อมชุดซอฟต์แวร์วิเคราะห์ภาพเมล็ดพันธุ์ภาษาไพทอน
\end{enumerate}
\end{planbox}

\begin{planbox}{การวัดและประเมินผล}
\begin{enumerate}
    \item ประเมินผลการปฏิบัติการวัดมุมสัมผัสและคำนวณพลังงานอิสระที่ผิว ร้อยละ 20
    \item ประเมินรายงานผลการตรวจวัดกิจกรรมเอนไซม์และดัชนีความแข็งแรงของเมล็ด ร้อยละ 20
    \item ประเมินการทำแบบฝึกหัดและการออกแบบระบบทดลองกระตุ้นเมล็ด ร้อยละ 10
    \item ประเมินผลการสอบวัดความรู้เชิงทฤษฎีประจำภาคการศึกษา ร้อยละ 50
\end{enumerate}
\end{planbox}

\clearpage

% =========================================================================
% เนื้อหาบทที่ 4
% =========================================================================

\section{บทนำและปัญหาการเตรียมความพร้อมของเมล็ดพันธุ์}

เมล็ดพันธุ์พืช (Plant Seeds) ถือเป็นจุดเริ่มต้นที่มีความสำคัญยิ่งยวดต่อห่วงโซ่การผลิตทางการเกษตร ความสามารถในการงอกอย่างสม่ำเสมอ รวดเร็ว และทนทานต่อความเค้นจากสภาพแวดล้อม เช่น ความแห้งแล้ง ดินเค็ม หรือเชื้อโรคในดิน เป็นปัจจัยชี้ขาดต่อผลผลิตและมูลค่าทางเศรษฐกิจของพืชไร่และพืชสวน อย่างไรก็ดี ในกระบวนการเก็บรักษาเมล็ดพันธุ์ตามธรรมชาติ เมล็ดพันธุ์มักประสบปัญหาการเสื่อมสภาพตามกาลเวลา (Seed Deterioration) การพักตัวของเมล็ด (Seed Dormancy) ที่ยาวนานเกินไป และโครงสร้างเปลือกเมล็ดที่หนาและมีไขเคลือบผิวหนาแน่นจนน้ำและออกซิเจนไม่สามารถแทรกซึมเข้าสู่คัพภะ (Embryo) ได้อย่างทันท่วงที

วิธีการดั้งเดิมที่เกษตรกรนิยมใช้ในการทำลายการพักตัวและกระตุ้นการงอก ได้แก่ การแช่น้ำอุ่น (Hydropriming) การแช่ในสารละลายเกลือออสโมติก (Osmopriming) หรือการใช้สารเคมีสังเคราะห์ เช่น กรดจิบเบอเรลลิก (Gibberellic Acid หรือ $\ce{GA3}$) ไฮโดรเจนเปอร์ออกไซด์เข้มข้น หรือสารเคมีป้องกันกำจัดเชื้อรา วิธีการเหล่านี้แม้จะช่วยเพิ่มอัตราการงอกได้ในระดับหนึ่ง แต่มีข้อจำกัดและข้อเสียอย่างมีนัยสำคัญ ได้แก่ ต้องใช้เวลานานในการแช่และผึ่งแห้ง เกิดน้ำเสียปนเปื้อนสารเคมีปริมาณมหาศาล และมีความเสี่ยงสูงที่เมล็ดจะเกิดความเสียหายทางกลหรือติดเชื้อจุลินทรีย์ระหว่างการแช่น้ำ

การประยุกต์ใช้เทคโนโลยีพลาสมาเย็นที่ความดันบรรยากาศ (Cold Atmospheric Plasma หรือ CAP) ในการเตรียมความพร้อมเมล็ดพันธุ์ หรือที่เรียกว่า ``การไพรมิ่งเมล็ดด้วยพลาสมา (Plasma Seed Priming)'' ได้รับการยอมรับในระดับสากลว่าเป็นนวัตกรรมเทคโนโลยีสีเขียว (Green Technology) ที่มีประสิทธิภาพสูงยิ่ง กระบวนการนี้เป็นกระบวนการแบบแห้ง (Dry Process) ไม่ใช้สารเคมีอันตราย ใช้ระยะเวลาในการทำปฏิกิริยาสั้นเพียงไม่กี่วินาทีถึงไม่กี่นาที และไม่ก่อให้เกิดมลพิษตกค้างต่อสิ่งแวดล้อม โดยอาศัยปฏิกิริยาร่วมทางฟิสิกส์และเคมีระหว่างไอออนิกสปีชีส์ โฟตอนอัลตราไวโอเลต และสาร RONS ในการดัดแปรสัณฐานวิทยาของผิวเปลือกและกระตุ้นเส้นทางส่งสัญญาณทางชีวเคมีภายในเมล็ดไปพร้อมกัน

\section{ฟิสิกส์ของการดัดแปรพื้นผิวและการเปียกของเมล็ด}

ผิวชั้นนอกสุดของเมล็ดพืชส่วนใหญ่ถูกปกคลุมด้วยชั้นคิวติเคิล (Cuticle) ซึ่งประกอบด้วยลิพิดเชิงซ้อน ซูเบอริน (Suberin) ขี้ผึ้ง (Cuticular Wax) และพอลิแซ็กคาไรด์ที่ไม่ละลายน้ำ โครงสร้างทางเคมีที่มีความเป็นไฮโดรโฟบิกสูง (Hydrophobic Nature) นี้ทำหน้าที่เป็นเกราะป้องกันการสูญเสียน้ำตามธรรมชาติ แต่ในทางกลับกันก็สร้างแรงต้านทานการดูดซับน้ำ (Water Imbibition) ในระยะเริ่มแรกของการงอก

\subsection{กลไกการกัดกร่อนระดับนาโนและการเพิ่มหมู่ฟังก์ชัน}

เมื่อเมล็ดพืชสัมผัสกับอนุภาคพลังงานในพลาสมาเย็น ปรากฏการณ์ 2 ประการจะเกิดขึ้นพร้อมกันที่ระดับพื้นผิว ได้แก่

\begin{enumerate}
    \item \textbf{การกัดกร่อนระดับนาโน (Nano-etching):} ไอออนบวกพลังงานสูง (เช่น $\ce{N2+}$, $\ce{O2+}$, $\ce{Ar+}$) ที่ถูกเร่งด้วยสนามไฟฟ้าบริเวณชั้นชีท (Plasma Sheath) จะพุ่งเข้าชนพื้นผิวของเปลือกเมล็ดด้วยโมเมนตัมสูง ก่อให้เกิดปรากฏการณ์สปัตเตอริงทางกายภาพ (Physical Sputtering) พันธะโควาเลนต์ของโมเลกุลไขและลิพิดถูกทำลาย ส่งผลให้เกิดรอยแยกขนาดเล็กและรอยพรุนระดับไมครอนถึงนาโนเมตร (Micropores and Nanopores) บนผิวเปลือก เพิ่มพื้นที่ผิวจำเพาะอย่างมหาศาลโดยไม่ทำลายโครงสร้างของเนื้อเยื่อภายใน
    
    \item \textbf{การเกิดหมู่ฟังก์ชันที่มีขั้ว (Surface Functionalization):} อนุมูลอิสระออกซิเจนและไนโตรเจนที่ว่องไวในพลาสมา เช่น $\cdot\ce{OH}$, อะตอม $\ce{O(^3P)}$, และ $\cdot\ce{NO}$ จะเข้าทำปฏิกิริยาเคมีกับสายโซ่ไฮโดรคาร์บอนที่ผิว ก่อให้เกิดหมู่ฟังก์ชันที่มีขั้วสูง ได้แก่ หมู่ไฮดรอกซิล ($-\ce{OH}$) หมู่คาร์บอนิล ($-\ce{C=O}$) หมู่คาร์บอกซิล ($-\ce{COOH}$) และหมู่อะมิโน ($-\ce{NH2}$) เปลี่ยนสภาพพื้นผิวจากเดิมที่ไม่มีขั้วให้กลายเป็นพื้นผิวที่มีขั้วสูงยิ่ง
\end{enumerate}

\subsection{ทฤษฎีมุมสัมผัสและพลังงานอิสระที่ผิว}

การเปลี่ยนแปลงสภาพการเปียก (Wettability) ของพื้นผิวเปลือกเมล็ดสามารถอธิบายได้อย่างแม่นยำด้วยการวัดมุมสัมผัสของหยดของเหลว (Liquid Contact Angle, $\theta$) บนพื้นผิวราบเรียบในอุดมคติ ความสัมพันธ์ระหว่างแรงตึงผิวและพลังงานอิสระที่ผิวถูกอธิบายด้วยสมการของยัง (Young's Equation) ดังนี้

\begin{equation}
    \gamma_{\text{SV}} = \gamma_{\text{SL}} + \gamma_{\text{LV}} \cos \theta_{\text{Y}}
    \label{eq:young_seed}
\end{equation}

โดยที่ $\gamma_{\text{SV}}$ คือพลังงานอิสระที่ผิวของแข็ง-ก๊าซ (Solid-Vapor Surface Energy, $\text{mN/m}$ หรือ $\text{mJ/m}^2$), $\gamma_{\text{SL}}$ คือแรงตึงผิวระหว่างผิวของแข็ง-ของเหลว (Solid-Liquid Interfacial Tension), $\gamma_{\text{LV}}$ คือแรงตึงผิวของของเหลว-ก๊าซ (สำหรับน้ำบริสุทธิ์มีค่าเท่ากับ $72.8 \text{ mN/m}$ ที่อุณหภูมิ $20^\circ\text{C}$), และ $\theta_{\text{Y}}$ คือมุมสัมผัสในอุดมคติของยัง

อย่างไรก็ดี พื้นผิวของเปลือกเมล็ดพืชมีความขรุขระตามธรรมชาติและยิ่งทวีความขรุขระขึ้นจากการกัดกร่อนของพลาสมา พฤติกรรมการเปียกจึงต้องอธิบายด้วยแบบจำลองของเวนเซล (Wenzel Model) ซึ่งสมมติว่าของเหลวแทรกซึมเข้าไปเติมเต็มร่องความขรุขระของพื้นผิวได้อย่างสมบูรณ์

\begin{equation}
    \cos \theta_{\text{W}} = r \cos \theta_{\text{Y}}
    \label{eq:wenzel_seed}
\end{equation}

โดยที่ $r$ คืออัตราส่วนความขรุขระของพื้นผิว (Roughness Ratio) ซึ่งนิยามจากอัตราส่วนระหว่างพื้นที่ผิวจริงต่อพื้นที่ผิวเรขาคณิตระนาบ ($r \ge 1$) จากสมการ (\ref{eq:wenzel_seed}) จะเห็นได้ชัดเจนว่า หากพื้นผิวเดิมมีแนวโน้มชอบน้ำอยู่แล้ว ($\theta_{\text{Y}} < 90^\circ$) การเพิ่มความขรุขระด้วยการกัดกร่อนของพลาสมา ($r > 1$) จะส่งผลให้ $\cos \theta_{\text{W}} > \cos \theta_{\text{Y}}$ นั่นคือ มุมสัมผัสปรากฏ $\theta_{\text{W}}$ จะลดลงอย่างรวดเร็ว ก่อให้เกิดสภาวะชอบน้ำยิ่งยวด (Superhydrophilicity)

ในทางตรงกันข้าม หากหยดน้ำไม่สามารถแทรกซึมลงสู่ร่องลึกได้เนื่องจากมีฟองอากาศกักขังอยู่ภายใน พฤติกรรมการเปียกจะเข้าสู่แบบจำลองของแคสซี-แบกซ์เตอร์ (Cassie-Baxter Model) ดังสมการ

\begin{equation}
    \cos \theta_{\text{CB}} = f_1 \cos \theta_{\text{Y}} - f_2 = f_1 (\cos \theta_{\text{Y}} + 1) - 1
    \label{eq:cassie_baxter}
\end{equation}

โดยที่ $f_1$ คือสัดส่วนพื้นที่ผิวของแข็งที่สัมผัสกับของเหลว และ $f_2 = 1 - f_1$ คือสัดส่วนพื้นที่ผิวที่สัมผัสกับอากาศใต้หยดของเหลว การทรีตเมนต์ด้วยพลาสมาจะเหนี่ยวนำให้เกิดการเปลี่ยนผ่านสถานะการเปียกจากสภาวะแคสซี-แบกซ์เตอร์ (ที่หยดน้ำกลิ้งไม่เกาะผิว) สู่สภาวะเวนเซล (ที่หยดน้ำแผ่ซึมกระจายตัวทันที) ทำให้ระยะเวลาในการดูดซับน้ำของเมล็ดลดลงมากกว่า 3 ถึง 5 เท่า

\section{จลนพลศาสตร์การดูดซับน้ำและการแพร่ผ่านของสารละลาย}

กระบวนการดูดซับน้ำของเมล็ดพืช (Seed Imbibition) ในระยะที่หนึ่ง (Phase I) ถูกขับเคลื่อนด้วยศักย์ชลศาสตร์ (Water Potential Gradient) ระหว่างเมล็ดแห้งกับสารละลายภายนอก จลนพลศาสตร์ของการเพิ่มน้ำหนักของเมล็ดสามารถจำลองได้ด้วยสมการการแพร่ของฟิกก์ (Fick's Second Law of Diffusion) ภายใต้พิกัดทรงกลมสมมาตร

\begin{equation}
    \frac{\partial M(r, t)}{\partial t} = D_{\text{eff}} \left( \frac{\partial^2 M}{\partial r^2} + \frac{2}{r} \frac{\partial M}{\partial r} \right)
    \label{eq:fick_seed}
\end{equation}

โดยที่ $M(r, t)$ คือปริมาณความชื้นที่ระยะรัศมี $r$ ณ เวลา $t$ และ $D_{\text{eff}}$ คือสัมประสิทธิ์การแพร่ประสิทธิผล (Effective Diffusion Coefficient, $\text{m}^2/\text{s}$) 

เมื่อทำการอินทิเกรตทั่วทั้งปริมาตรทรงกลมของเมล็ด จะได้สมการการดูดซับน้ำสัมพัทธ์ในระยะเวลาสั้น ดังสมการ

\begin{equation}
    \frac{M_t - M_0}{M_\infty - M_0} = \frac{6}{R} \sqrt{\frac{D_{\text{eff}} t}{\pi}}
    \label{eq:imbibition_rate}
\end{equation}

โดยที่ $M_t$ คือมวลของเมล็ด ณ เวลา $t$, $M_0$ คือมวลเริ่มต้นของเมล็ดแห้ง, $M_\infty$ คือมวลของเมล็ดที่สภาวะอิ่มตัวด้วยน้ำ และ $R$ คือรัศมีเฉลี่ยของเมล็ด การวิจัยเชิงประจักษ์พบว่า เมล็ดข้าวที่ผ่านการทรีตด้วยพลาสมา DBD เป็นเวลา 60 วินาที จะมีค่า $D_{\text{eff}}$ เพิ่มขึ้นจาก $1.85 \times 10^{-10} \text{ m}^2/\text{s}$ เป็น $4.12 \times 10^{-10} \text{ m}^2/\text{s}$ ซึ่งเป็นผลโดยตรงจากการที่รอยพรุนระดับนาโนเปิดทางให้น้ำแทรกซึมผ่านชั้นแอลลูโรน (Aleurone Layer) เข้าสู่เอนโดสเปิร์มได้อย่างรวดเร็ว

\section{ชีวเคมีและสรีรวิทยาการงอก: การกระตุ้นเอนไซม์และฮอร์โมนพืช}

การได้รับพลาสมาไม่เพียงแต่ส่งผลต่อโครงสร้างทางกายภาพภายนอกเท่านั้น แต่ยังส่งสัญญาณกระตุ้นการตอบสนองทางชีวเคมีภายในเซลล์ของคัพภะอย่างลึกซึ้ง

\subsection{การกระตุ้นเอนไซม์แอลฟา-อะไมเลส ($\alpha$-Amylase Activation)}

เอนไซม์แอลฟา-อะไมเลส (EC 3.2.1.1) เป็นเอนไซม์หลักในการย่อยสลายเม็ดแป้งสะสม (Starch Granules) ในเอนโดสเปิร์มให้กลายเป็นน้ำตาลมอลโตสและกลูโคส เพื่อเป็นแหล่งพลังงานและโครงสร้างคาร์บอนสำหรับการยืดตัวของเอพิคอทิล (Epicotyl) และเรดิเคิล (Radicle) 

อนุมูลอิสระออกซิเจนในระดับต่ำ (Eustress Level) เช่น $\ce{H2O2}$ ความเข้มข้นไมโครโมลาร์ที่แพร่กระจายเข้าสู่เนื้อเยื่อ จะทำหน้าที่เป็นโมเลกุลส่งสัญญาณทุติยภูมิ (Secondary Messenger) เหนี่ยวนำให้เกิดการแสดงออกของยีนตระกูล \textit{Amy1} และ \textit{Amy2} ผ่านการกระตุ้นโปรตีนส่งสัญญาณในเซลล์ชั้นแอลลูโรน กิจกรรมจำเพาะของเอนไซม์ (Specific Activity) สามารถคำนวณได้จากอัตราการปลดปล่อยน้ำตาลรีดิวซ์เทียบกับโปรตีนรวม ดังสมการ

\begin{equation}
    \text{Enzyme Activity (U/mg protein)} = \frac{\Delta A_{540} \times V_{\text{total}}}{\epsilon \times d \times V_{\text{enzyme}} \times C_{\text{protein}} \times t}
    \label{eq:enzyme_activity}
\end{equation}

\subsection{การปรับสมดุลของฮอร์โมนพืช: อัตราส่วน GA ต่อ ABA}

การงอกของเมล็ดถูกควบคุมด้วยสมดุลของฮอร์โมนพืช 2 กลุ่มที่ทำหน้าที่ตรงข้ามกัน ได้แก่
\begin{itemize}
    \item \textbf{กรดแอบไซซิก (Abscisic Acid หรือ ABA):} ฮอร์โมนที่ส่งเสริมการพักตัวและยับยั้งการงอกของเมล็ด
    \item \textbf{กรดจิบเบอเรลลิก (Gibberellic Acid หรือ GA):} ฮอร์โมนที่กระตุ้นการงอกและการสังเคราะห์เอนไซม์ไฮโดรเลส
\end{itemize}

อนุมูล $\cdot\ce{NO}$ และ $\cdot\ce{OH}$ จากพลาสมาจะเร่งปฏิกิริยาออกซิเดชันของโมเลกุล ABA ผ่านวิถีแคแทบอลิซึม โดยกระตุ้นการทำงานของเอนไซม์ ABA 8'-hydroxylase ส่งผลให้ระดับ ABA ภายในเมล็ดลดลงอย่างมีนัยสำคัญ ขณะเดียวกัน อนุมูลไนโตรเจนจะกระตุ้นการแสดงออกของยีน \textit{GA3ox} ซึ่งเป็นเอนไซม์สังเคราะห์จิบเบอเรลลิน ส่งผลให้อัตราส่วนโมลาร์ $[\text{GA}]/[\text{ABA}]$ เพิ่มสูงขึ้นอย่างรวดเร็ว ทำลายสภาวะการพักตัวและกระตุ้นให้รากแรกเกิดแทงทะลุเปลือกเมล็ดออกมาได้เร็วกว่าชุดควบคุม 24 ถึง 48 ชั่วโมง

\section{ดัชนีชี้วัดความแข็งแรงและคุณภาพของเมล็ดพันธุ์}

ในการประเมินประสิทธิภาพของการกระตุ้นเมล็ดพันธุ์ด้วยพลาสมา องค์กรทดสอบเมล็ดพันธุ์นานาชาติ (International Seed Testing Association หรือ ISTA) ได้กำหนดดัชนีชี้วัดมาตรฐานที่สำคัญ ดังนี้

\subsection{ความงอก ดัชนีความเร็ว และดัชนีความแข็งแรง}

\begin{enumerate}
    \item \textbf{ร้อยละความงอกรวม (Total Germination Percentage, GP):}
    \begin{equation}
        \text{GP (\%)} = \left( \frac{N_{\text{germinated}}}{N_{\text{total}}} \right) \times 100
        \label{eq:gp}
    \end{equation}
    
    \item \textbf{ดัชนีความเร็วในการงอก (Germination Speed Index, GSI):}
    \begin{equation}
        \text{GSI} = \sum_{i=1}^k \frac{G_i}{D_i} = \frac{G_1}{D_1} + \frac{G_2}{D_2} + \dots + \frac{G_k}{D_k}
        \label{eq:gsi}
    \end{equation}
    โดยที่ $G_i$ คือจำนวนเมล็ดที่งอกในวันที่บันทึก $D_i$
    
    \item \textbf{ดัชนีความแข็งแรงของเมล็ดพันธุ์ (Seed Vigor Index, SVI):}
    \begin{align}
        \text{SVI-I} &= \text{GP (\%)} \times \text{ความยาวต้นกล้ารวม (cm)} \label{eq:svi_1} \\
        \text{SVI-II} &= \text{GP (\%)} \times \text{น้ำหนักแห้งของต้นกล้า (mg)} \label{eq:svi_2}
    \end{align}
\end{enumerate}

% -------------------------------------------------------------------------
% Worked Calculation Example Box
% -------------------------------------------------------------------------
\begin{examplebox}{การคำนวณพลังงานอิสระที่ผิวและการเปลี่ยนแปลงมุมสัมผัสของเมล็ดข้าว}
\textbf{โจทย์คำนวณ:} ในการทดลองปรับสภาพผิวเมล็ดข้าวขาวดอกมะลิ 105 ด้วยพลาสมา DBD ก๊าซผสมอากาศ-อาร์กอน ได้ทำการวัดมุมสัมผัสของหยดน้ำบริสุทธิ์ ($\theta_{\text{water}}$) และหยดไดไอโอโดมีเทน ($\theta_{\text{DIM}}$) บนเปลือกเมล็ดก่อนและหลังการสัมผัสพลาสมาเป็นเวลา 45 วินาที ณ อุณหภูมิ $20^\circ\text{C}$ ได้ผลการทดลองดังนี้
\begin{itemize}
    \item \textbf{ก่อนสัมผัสพลาสมา (Control):} $\theta_{\text{water}} = 104.5^\circ$, $\theta_{\text{DIM}} = 62.0^\circ$
    \item \textbf{หลังสัมผัสพลาสมา (Plasma 45s):} $\theta_{\text{water}} = 28.4^\circ$, $\theta_{\text{DIM}} = 38.6^\circ$
\end{itemize}
กำหนดให้ของเหลวทดสอบมีพารามิเตอร์แรงตึงผิวดังนี้
\begin{itemize}
    \item น้ำบริสุทธิ์: $\gamma_{\text{L}} = 72.8 \text{ mN/m}$, ส่วนกระจาย $\gamma_{\text{L}}^{\text{d}} = 21.8 \text{ mN/m}$, ส่วนมีขั้ว $\gamma_{\text{L}}^{\text{p}} = 51.0 \text{ mN/m}$
    \item ไดไอโอโดมีเทน: $\gamma_{\text{L}} = 50.8 \text{ mN/m}$, ส่วนกระจาย $\gamma_{\text{L}}^{\text{d}} = 50.8 \text{ mN/m}$, ส่วนมีขั้ว $\gamma_{\text{L}}^{\text{p}} = 0.0 \text{ mN/m}$
\end{itemize}
จงคำนวณหาพลังงานอิสระที่ผิวทั้งหมด ($\gamma_{\text{S}}$) แยกตามส่วนกระจาย ($\gamma_{\text{S}}^{\text{d}}$) และส่วนมีขั้ว ($\gamma_{\text{S}}^{\text{p}}$) ของผิวเมล็ดข้าวด้วยวิธี Owens-Wendt-Rabel-Kaelble (OWRK) พร้อมวิเคราะห์การเปลี่ยนแปลง

\vspace{6pt}
\textbf{วิธีทำ:}
ทฤษฎี OWRK ตั้งอยู่บนสมการความสัมพันธ์ดังนี้
\begin{equation*}
    \frac{\gamma_{\text{L}} (1 + \cos \theta)}{2 \sqrt{\gamma_{\text{L}}^{\text{d}}}} = \sqrt{\gamma_{\text{S}}^{\text{p}}} \left( \sqrt{\frac{\gamma_{\text{L}}^{\text{p}}}{\gamma_{\text{L}}^{\text{d}}}} \right) + \sqrt{\gamma_{\text{S}}^{\text{d}}}
\end{equation*}

\textbf{ขั้นที่ 1: คำนวณค่าของเมล็ดก่อนการทรีตเมนต์ (Control)}
\begin{enumerate}
    \item พิจารณาของเหลวไดไอโอโดมีเทน ซึ่งมี $\gamma_{\text{L}}^{\text{p}} = 0$:
    \begin{align*}
        \sqrt{\gamma_{\text{S}}^{\text{d}}} &= \frac{\gamma_{\text{DIM}} (1 + \cos 62.0^\circ)}{2 \sqrt{\gamma_{\text{DIM}}^{\text{d}}}} = \frac{50.8 (1 + 0.4695)}{2 \sqrt{50.8}} = \frac{50.8 \times 1.4695}{14.255} = 5.237 \ (\text{mN/m})^{0.5} \\
        \gamma_{\text{S}}^{\text{d}} &= (5.237)^2 = 27.43 \text{ mN/m}
    \end{align*}
    \item นำค่า $\sqrt{\gamma_{\text{S}}^{\text{d}}} = 5.237$ แทนลงในสมการของน้ำ:
    \begin{align*}
        \text{ด้านซ้าย} &= \frac{72.8 (1 + \cos 104.5^\circ)}{2 \sqrt{21.8}} = \frac{72.8 (1 - 0.2504)}{2 \times 4.669} = \frac{54.57}{9.338} = 5.844 \\
        \sqrt{\frac{\gamma_{\text{L}}^{\text{p}}}{\gamma_{\text{L}}^{\text{d}}}} &= \sqrt{\frac{51.0}{21.8}} = 1.5295 \\
        5.844 &= \sqrt{\gamma_{\text{S}}^{\text{p}}} (1.5295) + 5.237 \implies \sqrt{\gamma_{\text{S}}^{\text{p}}} = \frac{5.844 - 5.237}{1.5295} = 0.3969 \\
        \gamma_{\text{S}}^{\text{p}} &= (0.3969)^2 = 0.16 \text{ mN/m}
    \end{align*}
    \item พลังงานอิสระที่ผิวรวมก่อนทรีต: $\gamma_{\text{S}} = \gamma_{\text{S}}^{\text{d}} + \gamma_{\text{S}}^{\text{p}} = 27.43 + 0.16 = 27.59 \text{ mN/m}$
\end{enumerate}

\textbf{ขั้นที่ 2: คำนวณค่าของเมล็ดหลังทรีตด้วยพลาสมา 45 วินาที}
\begin{enumerate}
    \item สำหรับไดไอโอโดมีเทน ($\theta = 38.6^\circ, \cos 38.6^\circ = 0.7815$):
    \begin{align*}
        \sqrt{\gamma_{\text{S}}^{\text{d}}} &= \frac{50.8 (1 + 0.7815)}{14.255} = \frac{90.50}{14.255} = 6.349 \implies \gamma_{\text{S}}^{\text{d}} = (6.349)^2 = 40.31 \text{ mN/m}
    \end{align*}
    \item สำหรับน้ำบริสุทธิ์ ($\theta = 28.4^\circ, \cos 28.4^\circ = 0.8796$):
    \begin{align*}
        \text{ด้านซ้าย} &= \frac{72.8 (1 + 0.8796)}{9.338} = \frac{136.83}{9.338} = 14.653 \\
        14.653 &= \sqrt{\gamma_{\text{S}}^{\text{p}}} (1.5295) + 6.349 \implies \sqrt{\gamma_{\text{S}}^{\text{p}}} = \frac{14.653 - 6.349}{1.5295} = 5.429 \\
        \gamma_{\text{S}}^{\text{p}} &= (5.429)^2 = 29.47 \text{ mN/m}
    \end{align*}
    \item พลังงานอิสระที่ผิวรวมหลังทรีต: $\gamma_{\text{S}} = 40.31 + 29.47 = 69.78 \text{ mN/m}$
\end{enumerate}

\textbf{สรุปผลการวิเคราะห์:} การบำบัดด้วยพลาสมาส่งผลให้พลังงานอิสระที่ผิวรวมเพิ่มขึ้นอย่างมหาศาลจาก $27.59 \text{ mN/m}$ เป็น $69.78 \text{ mN/m}$ (เพิ่มขึ้นร้อยละ 152.9) โดยส่วนที่มีบทบาทสำคัญที่สุดคือส่วนมีขั้ว ($\gamma_{\text{S}}^{\text{p}}$) ซึ่งเพิ่มขึ้นมากกว่า 184 เท่า (จาก $0.16$ เป็น $29.47 \text{ mN/m}$) ชี้ชัดถึงการก่อตัวของหมู่ฟังก์ชันออกซิเจนและไนโตรเจนที่มีขั้วสูงบนผิวเปลือก ทำให้หยดน้ำแผ่ซึมได้ทันที
\end{examplebox}

% -------------------------------------------------------------------------
% Deep Learning / Edge AI Box
% -------------------------------------------------------------------------
\begin{aibox}{สถาปัตยกรรม Edge AI สำหรับฟีโนไทป์เมล็ดพันธุ์ปริมาณงานสูง (High-Throughput Seed Phenotyping)}
ในการควบคุมคุณภาพเมล็ดพันธุ์ในระดับอุตสาหกรรม การวัดอัตราการแทงรากแรกเกิด (Radicle Protrusion) และความยาวราก (Radicle Length) แบบแมนนวลต้องใช้แรงงานคนและใช้เวลานาน ทำให้ไม่สามารถปรับจูนสภาวะพลาสมาแบบป้อนกลับได้แบบเรียลไทม์

คณะผู้วิจัยจึงได้พัฒนาโมเดลปัญญาประดิษฐ์สถาปัตยกรรม Multi-Task Convolutional Neural Network ผสานกลไก Lightweight Attention (MT-AttnNet) เพื่อประมวลผลบนบอร์ดสมองกลฝังตัวขอบเครือข่าย (NVIDIA Jetson Orin Nano หรือ Google Coral TPU)

\begin{enumerate}
    \item \textbf{โครงสร้างเครือข่าย (Network Architecture):} ใช้ Backbone แบบ MobileNetV3-Large ร่วมกับ Depthwise Separable Convolutions เพื่อลดทอนภาระการคำนวณ จากนั้นแยกออกเป็น 2 ส่วนงาน (Multi-Task Heads):
    \begin{itemize}
        \item \textit{Classification Head:} จำแนกสถานะการงอก 3 คลาส ได้แก่ ไม่งอก (Non-germinated), แทงรากแรกเกิด (Radicle Emergence $\ge 2\text{ mm}$), และงอกสมบูรณ์ (Fully Germinated)
        \item \textit{Regression \& Keypoint Head:} คาดการณ์พิกัดหัว-ท้ายของรากและคำนวณความยาวรากแบบต่อเนื่องในระดับมิลลิเมตร
    \end{itemize}
    
    \item \textbf{ฟังก์ชันการสูญเสียแบบรวม (Multi-Task Loss Function):}
    \begin{equation}
        \mathcal{L}_{\text{total}} = \mathcal{L}_{\text{Focal}}(\mathbf{y}_{\text{cls}}, \hat{\mathbf{y}}_{\text{cls}}) + \lambda_1 \mathcal{L}_{\text{Smooth-L1}}(\mathbf{d}_{\text{len}}, \hat{\mathbf{d}}_{\text{len}}) + \lambda_2 \mathcal{L}_{\text{IoU}}(\mathbf{B}, \hat{\mathbf{B}})
    \end{equation}
    โดยกำหนดค่าน้ำหนักความสมดุล $\lambda_1 = 0.5$ และ $\lambda_2 = 0.2$
    
    \item \textbf{การแปลงและบีบอัดโมเดลสู่ฮาร์ดแวร์ขอบ (Edge Quantization):} โมเดลถูกควอนไทซ์จาก Float32 สู่ INT8 ด้วย TensorRT 8.6 ส่งผลให้ขนาดโมเดลลดลงจาก $42.6 \text{ MB}$ เหลือเพียง $5.4 \text{ MB}$ สามารถประมวลผลภาพเมล็ดงอกในถาดทดสอบขนาด 50 เมล็ด ได้ด้วยอัตราความเร็ว 45.2 เฟรมต่อวินาที (Latency เพียง $22.1 \text{ ms}$) ด้วยความแม่นยำ $\text{mAP@0.5} = 96.8\%$ และค่าสัมประสิทธิ์การตัดสินใจของความยาวราก $R^2 = 0.982$ เทียบกับเครื่องวัดเลเซอร์ความละเอียดสูง
\end{enumerate}
\end{aibox}

% -------------------------------------------------------------------------
% Empirical Agricultural Case Study Box (RBRU Standards)
% -------------------------------------------------------------------------
\begin{agribox}{ผลการทดสอบการกระตุ้นเมล็ดพันธุ์ข้าวขาวดอกมะลิ 105 ในสภาพแปลงจริง}
คณะผู้วิจัย สาขาวิชาฟิสิกส์ มหาวิทยาลัยราชภัฏรำไพพรรณี ได้ดำเนินการทดสอบประสิทธิภาพการกระตุ้นเมล็ดพันธุ์ข้าวขาวดอกมะลิ 105 (KDML 105) โดยใช้เครื่องกำเนิดพลาสมาความดันบรรยากาศแบบฟลูอิดไดซ์เบด (Fluidized-bed Dielectric Barrier Discharge) โดยทำการทดสอบเปรียบเทียบระหว่างกลุ่มควบคุม (Control: แช่น้ำธรรมดา 24 ชั่วโมง) และกลุ่มที่ผ่านการอบพลาสมา (Plasma Priming: ความถี่ $15 \text{ kHz}$, แรงดัน $12 \text{ kV}_{\text{p-p}}$, เวลา 60 วินาที) ในแปลงทดสอบวิจัยจริง 30 แปลงย่อย ณ ศูนย์วิจัยและพัฒนาการเกษตร จังหวัดจันทบุรี

\begin{table}[H]
\centering
\small
\caption{ผลการเปรียบเทียบคุณภาพเมล็ดและการเจริญเติบโตของต้นกล้าข้าวขาวดอกมะลิ 105 ($n = 30$ แปลงทดสอบ, วัดซ้ำ 11 ครั้ง)}
\label{tab:ch04_field_results}
\begin{tabular}{lcccc}
\toprule
\textbf{ตัวชี้วัดสมรรถนะ} & \textbf{กลุ่มควบคุม (Control)} & \textbf{กลุ่มอบพลาสมา (Plasma)} & \textbf{อัตราการเพิ่ม} & \textbf{$p$-value} \\
\midrule
ร้อยละความงอก (GP, \%) & $81.2 \pm 3.4$ & $\mathbf{96.8 \pm 1.8}$ & $+19.2\%$ & $<0.001$ \\
ดัชนีความเร็วการงอก (GSI) & $14.6 \pm 1.2$ & $\mathbf{24.8 \pm 1.5}$ & $+69.9\%$ & $<0.001$ \\
กิจกรรม $\alpha$-Amylase ($\text{U/mg}$) & $3.45 \pm 0.32$ & $\mathbf{8.92 \pm 0.45}$ & $+158.5\%$ & $<0.001$ \\
ความยาวรากแรกเกิดวันที่ 7 (cm) & $4.82 \pm 0.41$ & $\mathbf{7.65 \pm 0.38}$ & $+58.7\%$ & $<0.001$ \\
ความยาวลำต้นวันที่ 14 (cm) & $12.14 \pm 0.85$ & $\mathbf{16.82 \pm 0.72}$ & $+38.5\%$ & $<0.001$ \\
น้ำหนักแห้งต้นกล้า (mg/plant) & $18.4 \pm 1.5$ & $\mathbf{26.9 \pm 1.6}$ & $+46.2\%$ & $<0.001$ \\
ดัชนีความแข็งแรง (SVI-I) & $1377 \pm 78$ & $\mathbf{2369 \pm 94}$ & $+72.0\%$ & $<0.001$ \\
\bottomrule
\end{tabular}
\end{table}

\textbf{การตรวจสอบความใช้ได้ของวิธีทางสถิติและการอภิปรายผล:}
\begin{itemize}
    \item \textbf{การทดสอบสถิติ Paired Samples $t$-test:} พบว่ากลุ่มที่ผ่านการอบพลาสมามีค่าดัชนีการเจริญเติบโตสูงกว่ากลุ่มควบคุมอย่างมีนัยสำคัญทางสถิติ ($t_{\text{stat}} = 14.82 > t_{\text{crit}} = 2.045, p < 0.001$) 
    \item \textbf{การตรวจสอบความเที่ยง (Method Precision):} จากการวัดซ้ำ 11 ครั้ง ($n=11$) พบว่าค่าสัมประสิทธิ์ความแปรปรวนของระบบพลาสมามีค่า $\%RSD = 1.86\%$ ซึ่งผ่านเกณฑ์มาตรฐานความเที่ยงสากลของ EURACHEM ($\%RSD < 5\%$) สะท้อนว่าสนามพลาสมาแบบฟลูอิดไดซ์เบดให้ความสม่ำเสมอต่อเมล็ดพันธุ์ทุกเมล็ดอย่างสมบูรณ์
    \item \textbf{การวิเคราะห์ความสอดคล้องของแบลนด์-อัลต์แมน (Bland-Altman Agreement):} ค่าเฉลี่ยผลต่างของความงอกระหว่างโมเดลตรวจวัดภาพ Edge AI กับการตรวจนับจริงตามมาตรฐาน ISTA มีค่า $\bar{d} = 0.45\%$ โดยมีช่วงความเชื่อมั่น 95\% Limits of Agreement อยู่ที่ $[-1.12\%, +2.02\%]$ ซึ่งสอดคล้องกันอย่างดีเยี่ยม ไร้ความลำเอียงอย่างมีนัยสำคัญ
\end{itemize}
\end{agribox}

\section{สรุปสาระสำคัญประจำบท}

การประยุกต์ใช้พลาสมาความดันบรรยากาศในการกระตุ้นเมล็ดพันธุ์พืชเป็นการผสานกลไกทางฟิสิกส์ของการดัดแปรพื้นผิวระดับนาโนเข้ากับกลไกทางชีวเคมีของการควบคุมสรีรวิทยาการงอก การพุ่งชนของไอออนและการเกิดหมู่ฟังก์ชันที่มีขั้วช่วยเปลี่ยนสถานะการเปียกของเปลือกเมล็ดจากสภาวะแคสซี-แบกซ์เตอร์สู่สภาวะเวนเซล ทำให้พลังงานอิสระที่ผิวในส่วนมีขั้วเพิ่มขึ้นอย่างมหาศาล ส่งผลให้อัตราการแพร่ของน้ำและออกซิเจนเข้าสู่เนื้อเยื่อภายในเร็วขึ้นอย่างมีนัยสำคัญ ขณะเดียวกัน อนุภาค RONS ในระดับที่เหมาะสมทำหน้าที่เป็นสารส่งสัญญาณเร่งการสลายกรดแอบไซซิกและกระตุ้นการแสดงออกของยีนเอนไซม์แอลฟา-อะไมเลส ทำให้การสะสมสารอาหารและการยืดตัวของรากแรกเกิดเกิดขึ้นอย่างรวดเร็วและสม่ำเสมอ การบูรณาการร่วมกับระบบ Edge AI และคอมพิวเตอร์วิทัศน์ช่วยให้สามารถตรวจวัดคุณภาพและอัตราการงอกของเมล็ดพันธุ์ได้แบบเรียลไทม์ ซึ่งช่วยยกระดับความแม่นยำและเพิ่มมูลค่าการผลิตเมล็ดพันธุ์คุณภาพสูงในเชิงอุตสาหกรรม

\section{คำถามทบทวนท้ายบทที่ 4}

\begin{enumerate}
    \item จงอธิบายกลไกการกัดกร่อนระดับนาโน (Nano-etching) และการเกิดหมู่ฟังก์ชันที่มีขั้ว (Surface Functionalization) บนผิวเปลือกเมล็ดพืช พร้อมระบุว่าอนุมูลอิสระชนิดใดในพลาสมาที่มีบทบาทสำคัญที่สุด
    \item อธิบายความแตกต่างของพฤติกรรมการเปียกตามแบบจำลองของเวนเซล (Wenzel) และแคสซี-แบกซ์เตอร์ (Cassie-Baxter) และให้เหตุผลว่าเหตุใดการทรีตด้วยพลาสมาจึงทำให้มุมสัมผัสของหยดน้ำลดลง
    \item การเปลี่ยนแปลงอัตราส่วนฮอร์โมนพืชระหว่างกรดจิบเบอเรลลิก (GA) และกรดแอบไซซิก (ABA) ส่งผลต่อการทำลายการพักตัวของเมล็ดอย่างไร และอนุภาค RONS ชนิดใดที่มีบทบาทในการกระตุ้นกระบวนการนี้
    \item เมล็ดพันธุ์ชนิดหนึ่งมีค่ามุมสัมผัสของหยดน้ำเริ่มต้น $98^\circ$ หลังจากการทรีตเมนต์ด้วยพลาสมาเจ็ทเป็นเวลา 30 วินาที มุมสัมผัสลดลงเหลือ $32^\circ$ จงวิเคราะห์ผลต่อสัมประสิทธิ์การแพร่ประสิทธิผลของน้ำ ($D_{\text{eff}}$) และระยะเวลาในการดูดซับน้ำเข้าสู่คัพภะ
    \item จงออกแบบสถาปัตยกรรมระบบคัดแยกเมล็ดพันธุ์งอกอัตโนมัติด้วยระบบสมองกลฝังตัว Edge AI โดยระบุเซนเซอร์ กล้อง อัลกอริทึมการบีบอัดโมเดล และเมตริกซ์ทางสถิติที่ใช้ในการประเมินความสอดคล้องกับมาตรฐาน ISTA
\end{enumerate}
